\subsubsection{Resolución del paso 3: del sistema de la cámara al del objeto}\label{sec:casco-ejes}

Los puntos $\vect{P}$ están en los ejes de la cámara ($X$ a la derecha, $Y$ hacia abajo, $Z$ hacia adelante). El visor 3D usa los ejes
del objeto ($x$ a la derecha, $y$ hacia arriba, $z$ hacia el observador). Se pasa de uno a otro con una traslación al centro del
objeto, $\vect{c}$ (el punto medio de la caja que contiene a todos los vértices), y una rotación de $180^\circ$ alrededor del eje $x$:
\begin{equation}
\vect{p}=R_o(\vect{P}-\vect{c}),\qquad R_o=\begin{pmatrix} 1 & 0 & 0 \\ 0 & -1 & 0 \\ 0 & 0 & -1 \end{pmatrix},\qquad \vect{c}=\begin{pmatrix} -3{,}354 \\ -36{,}101 \\ 379{,}535 \end{pmatrix}\ \text{mm}.
\end{equation}
$R_o$ es una rotación ($R_o^TR_o=I$ y $\det R_o=1$), su inversa es su transpuesta y, como es diagonal, $R_o^{-1}=R_o^T=R_o$. Para el píxel A:
\begin{equation}
\vect{p}_A=R_o\left(\begin{pmatrix} 3{,}195 \\ -25{,}557 \\ 300{,}871 \end{pmatrix}-\begin{pmatrix} -3{,}354 \\ -36{,}101 \\ 379{,}535 \end{pmatrix}\right)=R_o\begin{pmatrix} 6{,}549 \\ 10{,}544 \\ -78{,}664 \end{pmatrix}=\begin{pmatrix} 6{,}549 \\ -10{,}544 \\ 78{,}664 \end{pmatrix}\ \text{mm}.
\end{equation}

\paragraph{Con una sola matriz de $4\times4$.} La traslación y la rotación se pueden juntar en una matriz homogénea de $4\times4$ que actúa sobre
el punto con un $1$ agregado, $(\vect{P},1)$ (coordenadas homogéneas, igual que el píxel $(u,v,1)$):
\begin{equation}
\begin{pmatrix}\vect{p}\\1\end{pmatrix}=\mathcal{M}_o\begin{pmatrix}\vect{P}\\1\end{pmatrix},\qquad
\mathcal{M}_o=\begin{pmatrix}R_o&-R_o\vect{c}\\\vect{0}^{T}&1\end{pmatrix}=\begin{pmatrix} 1 & 0 & 0 & 3{,}354 \\ 0 & -1 & 0 & -36{,}101 \\ 0 & 0 & -1 & 379{,}535 \\ 0 & 0 & 0 & 1 \end{pmatrix} .
\end{equation}
Su determinante es $\det\mathcal{M}_o=\det R_o=1$ (una matriz de rotación y traslación no cambia volúmenes), y su inversa se obtiene por bloques,
porque $-R_o^{T}(-R_o\vect{c})=\vect{c}$:
\begin{equation}
\mathcal{M}_o^{-1}=\begin{pmatrix}R_o^{T}&\vect{c}\\\vect{0}^{T}&1\end{pmatrix}=\begin{pmatrix} 1 & 0 & 0 & -3{,}354 \\ 0 & -1 & 0 & -36{,}101 \\ 0 & 0 & -1 & 379{,}535 \\ 0 & 0 & 0 & 1 \end{pmatrix},\qquad
\mathcal{M}_o\,\mathcal{M}_o^{-1}=\begin{pmatrix} 1 & 0 & 0 & 0 \\ 0 & 1 & 0 & 0 \\ 0 & 0 & 1 & 0 \\ 0 & 0 & 0 & 1 \end{pmatrix}=I_4 .
\end{equation}
Con el píxel A, $\mathcal{M}_o(\vect{P}_A,1)=\begin{pmatrix} 6{,}549 \\ -10{,}544 \\ 78{,}664 \\ 1 \end{pmatrix}=(\vect{p}_A,1)$, y al revés, $\mathcal{M}_o^{-1}(\vect{p}_A,1)=\begin{pmatrix} 3{,}195 \\ -25{,}557 \\ 300{,}871 \\ 1 \end{pmatrix}=(\vect{P}_A,1)$: la inversa
devuelve el punto a los ejes de la cámara.

La figura~\ref{fig:casco-corte} muestra un corte lateral de los puntos del frente y de la parte de atrás y los rayos que los generan.

\begin{figure}[htbp]
\centering
\includegraphics[width=0.78\linewidth]{casco_corte.jpg}
\caption{Corte lateral (plano $Y$--$Z$) a lo largo de la columna $u=450$: cada punto del frente está sobre el rayo de su
píxel, a la profundidad $Z_{\text{frente}}$; la parte de atrás queda a $Z_{\text{atrás}}$.}
\label{fig:casco-corte}
\end{figure}
